\begin{frame}{왜 ``예시 몇 개''가 새 작업을 가르치는가}
  \begin{itemize}
    \item 예시만 포함시키면 새 형식의 작업을 수행하게
      된다는 사실(few-shot)은, \textbf{사전학습이 얼마나
      폭넓은 능력을 압축해뒀는가}의 증거.
    \item 사전학습 데이터에는 ``질문-답변'', ``예시-패턴''
      형태가 이미 \textbf{무수히} 포함 --- 프롬프트에
      비슷한 패턴을 보여주면 그 패턴을 이어가는 방향으로
      다음 토큰을 예측.
    \item 즉: \textbf{새 지식을 학습하는 것이 아니라}, 이미
      갖고 있는 능력 중 \textbf{어떤 것을 꺼내 쓸지}를
      프롬프트가 가리키는 것에 가까움.
    \item 17.1 관점: LLM 출력은 $P(x_t \mid x_1, \ldots,
      x_{t-1})$ --- \textbf{지금까지 프롬프트에 포함된
      모든 토큰}에 대한 조건부 확률. 예시 3개를 붙이면
      조건이 \textbf{그 예시를 포함한 전체 시퀀스}.
    \item Week 12.2 어텐션 관점: 질문 토큰이 다음 단어를
      고를 때 예시 부분을 \textbf{실제로 ``살펴보는''} 것
      --- 그래서 예시를 바꾸면 예측이 바뀜.
    \item ``few-shot = 모델이 예시로 \textbf{학습}한다''는
      표현은 \textbf{정확하지 않다} --- 컨텍스트 안의 토큰이
      \textbf{조건}으로 작용할 뿐, 파라미터에 새겨지는 것이
      아님. 프롬프트를 바꾸면 \textbf{즉시, 완벽하게}
      원래대로 복귀.
  \end{itemize}
  \vskip0.3em
  {\footnotesize 조건으로 삼는 토큰이 늘면 어텐션 계산량이
  $n^2$으로 커진다는 비용 함의도 따라옴(Week 12.3).}
\end{frame}
